\documentclass[a4paper,11pt]{article}
\usepackage{iftex}
\usepackage{amssymb}
\ifLuaTeX
  \providecommand{\DeclareUnicodeCharacter}[2]{}
\fi
\DeclareUnicodeCharacter{202F}{\,}% NARROW NO-BREAK SPACE (z. B. in „z. B.“)
\DeclareUnicodeCharacter{2011}{-}% NON-BREAKING HYPHEN (geschützter Bindestrich)

\ifLuaTeX
  \usepackage{fontspec}
  \defaultfontfeatures{Ligatures=TeX}
  \usepackage[ngerman]{babel}
\else
  \usepackage[T1]{fontenc}
  \usepackage[utf8]{inputenc}
  \usepackage[ngerman]{babel}
\fi
\usepackage{csquotes}
\usepackage{geometry}\geometry{margin=2.5cm}
\usepackage[final]{microtype}
\emergencystretch=1.5em
\usepackage{xurl}\urlstyle{same}
\usepackage[hidelinks,bookmarksopen=true,bookmarksnumbered=true]{hyperref}
\usepackage{cleveref}
\usepackage{enumitem}
\usepackage[backend=biber,style=ieee,sorting=nyt,maxnames=6]{biblatex}
\usepackage{newunicodechar}

\addbibresource{bib/seminar.bib}

% Zentrale Einstellungen fuer den Themenkatalog
% Diese Datei ist der erste Anlaufpunkt fuer semesterweise Anpassungen.

\newcommand{\CatalogSemester}{SoSe 2026}
\newcommand{\CatalogVersion}{v2.0}
\newcommand{\CatalogMaintainer}{Dozententeam Ingenieur-Informatik}
\newcommand{\CatalogContactEmail}{michael.kreutzer@mni.thm.de}

% false = Studierendenansicht (nur aktive Themen)
% true  = Dozierendenansicht (auch hidden/used/draft sichtbar)
\newif\ifTeacherView
\TeacherViewfalse


\hypersetup{
  pdftitle={CS1025 Hauptseminar -- Themenkatalog (\CatalogSemester)},
  pdfauthor={Ingenieur-Informatik \textperiodcentered{} THM},
  pdfsubject={Modular gepflegter Seminarkatalog mit Registry- und Statuslogik},
  pdfkeywords={Themenkatalog, Seminar, Informatik, Ingenieur-Informatik, THM, \CatalogSemester}
}

\usepackage{xparse}
\usepackage{pgfkeys}
\usepackage{etoolbox}
\usepackage{xcolor}

\pgfkeys{/topic/.is family, /topic,
  id/.initial = ,
  difficulty/.initial = ,
  type/.initial = ,
  prerequisites/.initial = ,
  practice/.initial = ,
  owner/.initial = ,
  updated/.initial = ,
  tags/.initial = ,
}

\newcommand{\topicset}[1]{\pgfkeys{/topic,#1}}
\newcommand{\topicgetval}[2]{\pgfkeysgetvalue{/topic/#1}{#2}\ifstrempty{#2}{\def#2{--}}{}}

\newif\ifTopicDesc \newif\ifTopicKW \newif\ifTopicQ \newif\ifTopicScope \newif\ifTopicPractice
\newcommand{\CurrentTopicStatus}{active}
\newcommand{\SetCurrentTopicStatus}[1]{\gdef\CurrentTopicStatus{#1}}

\newcommand{\TopicDescription}[1]{\TopicDesctrue\paragraph{Beschreibung}#1}
\newcommand{\TopicKeywords}[1]{\TopicKWtrue\paragraph{Schlagwörter}#1}
\newcommand{\TopicLeitfragen}[1]{\TopicQtrue\paragraph{Leitfragen}\begin{itemize}#1\end{itemize}}
\newcommand{\TopicScope}[1]{\TopicScopetrue\paragraph{Abgrenzung (nicht enthalten)}\begin{itemize}#1\end{itemize}}
\newcommand{\TopicPractice}[1]{\TopicPracticetrue\paragraph{Optionale Praxisidee}#1}

\newcommand{\TopicStatusLabel}{%
  \ifTeacherView
    \ifstrequal{\CurrentTopicStatus}{active}{\textcolor{gray}{Status: aktiv}}{%
      \textcolor{red!70!black}{Status: \CurrentTopicStatus}%
    }%
  \fi
}

\newcommand{\TopicMeta}{%
  \begingroup
  \def\Tdifficulty{}\def\Ttype{}\def\Tprereq{}\def\Tpractice{}\def\Towner{}\def\Tupdated{}\def\Tid{}
  \topicgetval{difficulty}{\Tdifficulty}
  \topicgetval{type}{\Ttype}
  \topicgetval{prerequisites}{\Tprereq}
  \topicgetval{practice}{\Tpractice}
  \topicgetval{owner}{\Towner}
  \topicgetval{updated}{\Tupdated}
  \topicgetval{id}{\Tid}
  \noindent\textit{Niveau: \Tdifficulty \quad Typ: \Ttype \quad Praxisanteil: \Tpractice \quad Betreuer: \Towner \quad Aktualisiert: \Tupdated}
  \ifTeacherView\quad \TopicStatusLabel\fi
  \par
  \smallskip
  \noindent\textbf{Empfohlene Vorkenntnisse:} \Tprereq\par\medskip
  \endgroup
}

\newcommand{\TopicWarnIfMissingMeta}[1]{%
  \begingroup
    \def\Ttmp{}\topicgetval{#1}{\Ttmp}%
    \ifstrequal{\Ttmp}{--}{%
      \def\Tid{}\topicgetval{id}{\Tid}%
      \PackageWarning{topic}{Thema (id=\Tid): Metadatum '#1' fehlt!}%
    }{}%
  \endgroup
}

\NewDocumentEnvironment{topic}{ O{} m }{%
  \topicset{#1}%
  \TopicDescfalse \TopicKWfalse \TopicQfalse \TopicScopefalse \TopicPracticefalse%
  \subsection{#2}%
  \TopicMeta
  \begin{refsection}%
}{%
  \begingroup
    \def\Tid{}\topicgetval{id}{\Tid}%
    \ifTopicDesc\else\PackageWarning{topic}{Thema (id=\Tid): 'Beschreibung' fehlt!}\fi
    \ifTopicQ\else\PackageWarning{topic}{Thema (id=\Tid): 'Leitfragen' fehlen!}\fi
    \ifTopicScope\else\PackageWarning{topic}{Thema (id=\Tid): 'Abgrenzung' fehlt!}\fi
    \ifTopicPractice\else\PackageWarning{topic}{Thema (id=\Tid): 'Optionale Praxisidee' fehlt!}\fi
    \ifTopicKW\else\PackageWarning{topic}{Thema (id=\Tid): 'Schlagwörter' fehlen!}\fi
    \TopicWarnIfMissingMeta{difficulty}%
    \TopicWarnIfMissingMeta{type}%
    \TopicWarnIfMissingMeta{prerequisites}%
    \TopicWarnIfMissingMeta{practice}%
  \endgroup
  \printbibliography[heading=subbibliography,title={Literatur (Auswahl)}]%
  \end{refsection}%
  \medskip\hrule\medskip
}

\newcommand{\TopicEntry}[2]{%
  \SetCurrentTopicStatus{#1}%
  \ifstrequal{#1}{active}{%
    \input{#2}%
  }{%
    \ifTeacherView
      \input{#2}%
    \fi
  }%
  \SetCurrentTopicStatus{active}%
}

\setlist[itemize]{leftmargin=*,nosep}
\setlist[enumerate]{leftmargin=*,nosep}


\title{CS1025 Hauptseminar -- Themenkatalog (\CatalogSemester)}
\author{Ingenieur-Informatik \textperiodcentered{} THM}
\date{Stand: \today}

\begin{document}

\maketitle

\noindent\textbf{Hinweis zur Nutzung.}
Dieser Katalog zeigt in der Studierendenansicht nur aktuell freigegebene Themen.
Die interne Pflege erfolgt über statusgesteuerte Registry-Dateien, sodass Themen semesterweise
schnell aktiviert, ausgeblendet oder als bereits verwendet markiert werden können.
Bei inhaltlichen Rückfragen oder Errata bitte an
\href{mailto:\CatalogContactEmail}{\CatalogContactEmail} wenden.

\vspace{0.75\baselineskip}
\noindent\textbf{Katalogversion.}
Fassung: \texttt{\CatalogVersion} (\CatalogSemester)

\section*{Hauptseminar CS1025 -- Rahmen, Lernziele und Hinweise zur Themenwahl}
\subsection*{Lernziele}
Nach erfolgreicher Teilnahme können Studierende:
\begin{itemize}
  \item aktuelle technische Themen systematisch recherchieren und die Qualität ihrer Quellen bewerten,
  \item einen technischen Gegenstand klar abgrenzen, geeignete Leitfragen formulieren und daraus eine stringente Argumentation ableiten,
  \item wesentliche Sachverhalte zielgruppenorientiert präsentieren und in der Diskussion fundiert vertreten,
  \item Ergebnisse schriftlich nach wissenschaftlichen Standards dokumentieren und kritisch einordnen.
\end{itemize}

\subsection*{Erwartete Leistungen im Seminar}
Das Seminar umfasst typischerweise:
\begin{itemize}
  \item eine Präsentation von ca.\ 30 Minuten,
  \item eine anschließende Diskussion bzw. Fragerunde,
  \item eine schriftliche Ausarbeitung mit ca.\ 4500 Wörtern im Hauptteil,
  \item nachvollziehbare Quellenarbeit, saubere Zitation und eine klare Abgrenzung des Themas.
\end{itemize}

\subsection*{Hinweise zur Themenwahl}
Wähle ein Thema so, dass es fachlich interessant, aber innerhalb des Semesters realistisch bearbeitbar ist.
Hilfreiche Auswahlkriterien sind insbesondere:
\begin{itemize}
  \item ausreichende und belastbare Quellenlage,
  \item ein klar formulierbarer Fokus statt eines zu breiten Überblicks,
  \item ein erkennbarer Bezug zur Ingenieur-Informatik,
  \item ein sinnvoller Ergebnistyp, etwa Vergleich, Einordnung, Systemanalyse oder praxisnahe Bewertung.
\end{itemize}

\noindent\textit{Hinweis:} Die Metadaten jeder Themenkarte (Niveau, Typ, empfohlene Vorkenntnisse, Praxisanteil) dienen als Auswahlhilfe. Nicht jedes Thema erfordert einen praktischen Aufbau; entscheidend sind fachliche Einordnung, saubere Struktur und wissenschaftliche Nachvollziehbarkeit.


\tableofcontents
\clearpage

\section{Eingebettete Intelligenz und Edge Computing}
% Statuswerte: active | hidden | used | draft
% In der Studierendenansicht werden nur active-Themen angezeigt.

\TopicEntry{active}{tex/topics/embedded/embedded_vision_ondevice.tex}
\TopicEntry{active}{tex/topics/embedded/mlperf_tiny_bench.tex}
\TopicEntry{active}{tex/topics/embedded/ondevice_genai_edge.tex}
\TopicEntry{active}{tex/topics/embedded/pytorch_embedded.tex}
\TopicEntry{active}{tex/topics/embedded/riscv_auto_embedded.tex}
\TopicEntry{active}{tex/topics/embedded/sdv_zonal_ee_arch.tex}
\TopicEntry{active}{tex/topics/embedded/tensorflow_intro.tex}
\TopicEntry{active}{tex/topics/embedded/tinyml_tflm.tex}
\TopicEntry{hidden}{tex/topics/embedded/wio_terminal_tinyml.tex}


\section{Sichere und resiliente Systeme}
% Statuswerte: active | hidden | used | draft
% In der Studierendenansicht werden nur active-Themen angezeigt.

\TopicEntry{active}{tex/topics/security/cra_embedded_iot.tex}
\TopicEntry{active}{tex/topics/security/ics_energy_attacks.tex}
\TopicEntry{active}{tex/topics/security/iot_cyberattacks.tex}
\TopicEntry{active}{tex/topics/security/memfault_remote_debug.tex}
\TopicEntry{active}{tex/topics/security/ota_strategies.tex}
\TopicEntry{active}{tex/topics/security/safety_vs_security.tex}
\TopicEntry{active}{tex/topics/security/sidechannel_attacks.tex}


\section{Vernetzte Systeme und industrielles IoT}
% Statuswerte: active | hidden | used | draft
% In der Studierendenansicht werden nur active-Themen angezeigt.

\TopicEntry{active}{tex/topics/iot/bluetooth_mesh.tex}
\TopicEntry{active}{tex/topics/iot/canopen.tex}
\TopicEntry{active}{tex/topics/iot/dds_middleware.tex}
\TopicEntry{active}{tex/topics/iot/fiveg_m2m.tex}
\TopicEntry{active}{tex/topics/iot/lorawan_ttn.tex}
\TopicEntry{active}{tex/topics/iot/mqtt_pubsub.tex}
\TopicEntry{active}{tex/topics/iot/node_red.tex}
\TopicEntry{active}{tex/topics/iot/opcua.tex}
\TopicEntry{active}{tex/topics/iot/opcuafx_tsn_field.tex}
\TopicEntry{active}{tex/topics/iot/profinet.tex}


\section{Robotik, Assistenzsysteme und Autonomie}
% Statuswerte: active | hidden | used | draft
% In der Studierendenansicht werden nur active-Themen angezeigt.

\TopicEntry{active}{tex/topics/robotics/eta_aids.tex}
\TopicEntry{active}{tex/topics/robotics/indoor_nav.tex}
\TopicEntry{active}{tex/topics/robotics/logistics_drones.tex}
\TopicEntry{active}{tex/topics/robotics/ros2_micro.tex}
\TopicEntry{active}{tex/topics/robotics/ros_industrial.tex}
\TopicEntry{active}{tex/topics/robotics/service_robots.tex}


\section{Zukunftstechnologien und nachhaltige Systeme}
% Statuswerte: active | hidden | used | draft
% In der Studierendenansicht werden nur active-Themen angezeigt.

\TopicEntry{active}{tex/topics/future/ar_vr_xr.tex}
\TopicEntry{active}{tex/topics/future/neuroprostheses.tex}
\TopicEntry{active}{tex/topics/future/printing_3d.tex}
\TopicEntry{active}{tex/topics/future/quantum_intro.tex}
\TopicEntry{active}{tex/topics/future/starlink_megaconstellations.tex}


\section{Offene Themen und Studierendenprojekte}
% Statuswerte: active | hidden | used | draft
% In der Studierendenansicht werden nur active-Themen angezeigt.

\TopicEntry{active}{tex/topics/open/open_guided.tex}
\TopicEntry{active}{tex/topics/open/indoor_context.tex}
\TopicEntry{active}{tex/topics/open/tasmota_building.tex}



\end{document}
